\documentclass{article}
\usepackage{amsmath}
\usepackage{amsfonts}
\usepackage{hyperref}
\usepackage{graphicx}

\title{Introduction to Computing for Physicists}
\author{Alessandro Rayan Ahmed}

\begin{document}
\maketitle
\begin{abstract}
  A brief overview of the general ideas of computing. A practical guide for
  learning a language efficiently, focusing on the \textit{C} programming
  language. A sweep of programming knowledge to improve algorithms.
\end{abstract}


\begin{itemize}
\item An \textbf{operation} is involved for the change of \textit{data/information}.
\item An \textbf{algorithm} is a \textit{static} set of rules.
\item A \textbf{process} is the \textit{dynamics} emerging from applying the set of rules.
\end{itemize}

\section{Flux diagrams}
The flux-diagram is a convenient tool for \textit{designing} an
algorithm because the flow of the program is made explicit with
the directions of the arrows pointing.

\section{Positional numbering system}
\subsection{Integers}
The value of each digit depends on its position.
\begin{itemize}
\item $b$ represents the \textit{basis} of the system.
\item $\Gamma$ is the system's \textit{alphabet}, it must contain $b$ number of symbols.
\item $d_i \in \Gamma$ represents the digit in the i-th \textit{position}.
\end{itemize}

Big endian vs. small endian
\begin{equation}
  \label{eq:positional-representation}
  \begin{gathered}
    (D_nD_{n-1}\cdots D_0)_b \equiv \sum_{i=1}^n D_ib^i = D_0 \cdot b^0 + D_1 \cdot b^1 + \dots + D_n \cdot b^n \\
    (d_0d_1\cdots d_n)_b \equiv \sum_{i=1}^n d_ib^i = d_0 \cdot b^0 + d_1 \cdot b^1 + \dots + d_n \cdot b^n
  \end{gathered}
\end{equation}

\begin{itemize}
\item $\Gamma_{\text{hex}} = (0, 1, 2, 3, 4, 5, 6, 7, 8, 9, \text{a}, \text{b}, \text{c}, \text{d}, \text{e}, \text{f})$
\item $\Gamma_{\text{dec}} = (0, 1, 2, 3, 4, 5, 6, 7, 8, 9)$
\item $\Gamma_{\text{oct}} = (0, 1, 2, 3, 4, 5, 6, 7)$
\item $\Gamma_{\text{bin}} = (0, 1)$
\end{itemize}

\begin{center}
  \begin{tabular}{|r|r|r|r|}
    \hline
    Hexadecimal & Decimal & Octal & Binary \\
    \hline
    $0_{\text{hex}}$ & $0_{\text{dec}}$ & $0_{\text{oct}}$ & $0_{\text{bin}}$ \\
    $1_{\text{hex}}$ & $1_{\text{dec}}$ & $1_{\text{oct}}$ & $1_{\text{bin}}$ \\
    $2_{\text{hex}}$ & $2_{\text{dec}}$ & $2_{\text{oct}}$ & $10_{\text{bin}}$ \\
    $3_{\text{hex}}$ & $3_{\text{dec}}$ & $3_{\text{oct}}$ & $11_{\text{bin}}$ \\
    $4_{\text{hex}}$ & $4_{\text{dec}}$ & $4_{\text{oct}}$ & $100_{\text{bin}}$ \\
    $5_{\text{hex}}$ & $5_{\text{dec}}$ & $5_{\text{oct}}$ & $101_{\text{bin}}$ \\
    $6_{\text{hex}}$ & $6_{\text{dec}}$ & $6_{\text{oct}}$ & $110_{\text{bin}}$ \\
    $7_{\text{hex}}$ & $7_{\text{dec}}$ & $7_{\text{oct}}$ & $111_{\text{bin}}$ \\
    $8_{\text{hex}}$ & $8_{\text{dec}}$ & $10_{\text{oct}}$ & $1000_{\text{bin}}$ \\
    $9_{\text{hex}}$ & $9_{\text{dec}}$ & $11_{\text{oct}}$ & $1001_{\text{bin}}$ \\
    $\text{a}_{\text{hex}}$ & $10_{\text{dec}}$ & $12_{\text{oct}}$ & $1010_{\text{bin}}$ \\
    $\text{b}_{\text{hex}}$ & $11_{\text{dec}}$ & $13_{\text{oct}}$ & $1011_{\text{bin}}$ \\
    $\text{c}_{\text{hex}}$ & $12_{\text{dec}}$ & $14_{\text{oct}}$ & $1100_{\text{bin}}$ \\
    $\text{d}_{\text{hex}}$ & $13_{\text{dec}}$ & $15_{\text{oct}}$ & $1101_{\text{bin}}$ \\
    $\text{e}_{\text{hex}}$ & $14_{\text{dec}}$ & $16_{\text{oct}}$ & $1110_{\text{bin}}$ \\
    $\text{f}_{\text{hex}}$ & $15_{\text{dec}}$ & $17_{\text{oct}}$ & $1111_{\text{bin}}$ \\
    \hline
  \end{tabular}
\end{center}

\subsection{Rationals}
\paragraph{Fixed-point representation}
\paragraph{Mobile-point representation}
\paragraph{IEEE 754}
Sign mantissa significand

\section{Turing machines}

\end{document}
